% reproducibility.tex
\label{sec:reproducibility}

\paragraph{Environment.}
Python (with \texttt{sympy} and \texttt{numpy}); the engine cross-check additionally
imports \texttt{dcl\_core} (dcl-core v0.3.0). No GPU is required---the derivations
are symbolic and the extraction runs on a small lattice.

\paragraph{Verify the headline result.}
From a fresh clone, the central claims are reproduced by
%
\begin{quote}\ttfamily
python -u src/utilities/electric\_induced\_action.py
\end{quote}
%
which prints the magnetic anchor $\{4,4,16\}$, the electric block $\{1,4,4\}$ (axis
suppressed), $\mu^{-1}=\operatorname{adj}(\varepsilon)$, and the null birefringence
split, all \texttt{PASS}. The engine cross-check that $\varepsilon=P$ is what the
tick rule actually produces is
%
\begin{quote}\ttfamily
python -u src/experiments/exp\_01\_electric\_permittivity\_extraction.py
\end{quote}

\paragraph{Build the paper.}
From \texttt{paper/}: \texttt{pdflatex main} $\to$ \texttt{bibtex main} $\to$
\texttt{pdflatex main} (twice). The \verb|\input| equation fragments under
\texttt{sections/generated/} are regenerated by
\texttt{electric\_induced\_action.py}.

\paragraph{Where to look next.}
The open items and the momentum-space route are in
\texttt{src/utilities/trlnT\_prescription.py} and
\texttt{ward\_safe\_diagnosis.py} (with the matching notes); the independent
numerical confirmation of the birefringence verdict is in the companion
Paper~IV~\cite{PaperIV}.
